% Non-orthogonal basis projection (general formulation)
% Compile with: xelatex nonorthogonal_basis_projection.tex
\documentclass{article}
\usepackage[UTF8]{ctex}
% Overleaf (Linux) 環境用 Noto CJK TC，Windows 在地編譯可改成 PMingLiU / Microsoft JhengHei
\setCJKmainfont{Noto Serif CJK TC}
% \setCJKmainfont{Noto Sans CJK TC}     % 改用無襯線（黑體）
% \setCJKmainfont{PMingLiU}             % Windows 在地：新細明體
% \setCJKmainfont{Microsoft JhengHei}   % Windows 在地：微軟正黑體
\usepackage{amsmath}
\usepackage{amssymb}
\usepackage{bm}

\begin{document}

\subsubsection*{非正交基底投影}

假設存在一組非正交座標，將向量 $\mathbf{v}$ 分解到非正交基底向量
$\{\mathbf{e}_1, \mathbf{e}_2, \mathbf{e}_3\}$ 上：
\begin{equation}
\mathbf{v} \;=\; \lambda_1 \mathbf{e}_1 + \lambda_2 \mathbf{e}_2 + \lambda_3 \mathbf{e}_3
\qquad \text{（三維）} .
\end{equation}

將基底向量以 column 形式排列成矩陣，可改寫為
\begin{equation}
\mathbf{v}
\;=\;
\begin{bmatrix} \mathbf{e}_1 & \mathbf{e}_2 & \mathbf{e}_3 \end{bmatrix}
\begin{bmatrix} \lambda_1 \\ \lambda_2 \\ \lambda_3 \end{bmatrix} .
\end{equation}

兩邊左乘 $\bigl[\mathbf{e}_1\; \mathbf{e}_2\; \mathbf{e}_3\bigr]^{-1}$，
即可解出座標
\begin{equation}
\begin{bmatrix} \mathbf{e}_1 & \mathbf{e}_2 & \mathbf{e}_3 \end{bmatrix}^{-1}
\mathbf{v}
\;=\;
\begin{bmatrix} \lambda_1 \\ \lambda_2 \\ \lambda_3 \end{bmatrix} .
\end{equation}

\paragraph{通用形式（二維與三維皆適用）。}
令 $E = \bigl[\mathbf{e}_1\; \mathbf{e}_2\; \cdots\; \mathbf{e}_n\bigr]$ 為將
基底向量依 column 排列所形成之矩陣，
$\bm{\lambda} = (\lambda_1,\ldots,\lambda_n)^{\!\top}$ 為對應之座標向量，則
\begin{equation}
\mathbf{v} \;=\; E\,\bm{\lambda},
\qquad
\bm{\lambda} \;=\; E^{-1} \mathbf{v}
\quad\text{（$E$ 為方陣且基底線性獨立時）} .
\end{equation}
二維情形只需令 $E = \bigl[\mathbf{e}_1\; \mathbf{e}_2\bigr] \in \mathbb{R}^{2\times 2}$，
形式完全相同。

\paragraph{過定基底（over-determined）。}
當基底向量數量超過所嵌入空間的維度時（例如四支磁極軸投影至 2D 平面），
$E$ 不再為方陣，$E^{-1}$ 不存在。此時最小平方意義下的座標解為
\begin{equation}
\bm{\lambda} \;=\; (E^{\!\top} E)^{-1} E^{\!\top} \mathbf{v}
\;=\; E^{+}\,\mathbf{v} ,
\end{equation}
其中 $E^{+}$ 為 Moore--Penrose 偽逆。

\end{document}
